\section{The Model \label{sec:themodel}}

Dynamic general equilibrium models can explain many features of business cycles (\cite{Prescott:1986aa}). In \cites{Kydland:1982aa} closed-economy model, fluctuations arise from a single aggregate productivity shock. The model accounts for the magnitude of fluctuations in output, consumption and investment, as well as the correlations of consumption and investment with output. It reproduces two central empirical regularities: consumption and investment are procyclical, and investment is more volatile than output while consumption is less volatile.

\subsection{BKK's Expanded Model}
\cite{Backus:1992aa} extend this framework to two countries with imperfectly correlated technology shocks. Their aim is to explain international business cycles within the RBC framework. The expanded model must account for both the fluctuations and comovements within each country and the comovements across countries.

BKK (1993) use a simpler version of the two-country model developed in BKK (1992). The countries produce a single good and share the same preferences and production technology. Country-specific shocks to total factor productivity are the sole source of uncertainty and drive business cycle fluctuations. Households maximize lifetime utility, which depends on consumption and leisure, and supply labor. Firms combine labor and capital to produce output.

\Cref{tab:modelresult} reports simulations from four models. The first row presents \cites{Kydland:1982aa} closed-economy model. The second presents BKK's (1993) baseline two-country model, with complete markets and no transport costs. The remaining rows introduce transport costs and eliminate trade between the countries, respectively.

Access to international borrowing changes the behavior of consumption and investment because expenditure is no longer constrained by domestic production. Both variables become less correlated with output than in the closed-economy model. Investment also becomes more volatile: when resources can move freely across borders, the country experiencing the more favorable technology shock attracts more capital than it would in a closed economy.


\begin{table}[ht]
\centering
\begin{subtable}{\textwidth}
  \centering
  \caption{Standard Deviations and Within-Country Comovements}
  \begin{adjustbox}{max width=\textwidth}
  \begin{tabular}{lcccccccccccc}
  \toprule
          &       &       & \multicolumn{4}{c}{Ratio of $\sigma$ to  $\sigma_y$} & \multicolumn{1}{c}{} & \multicolumn{5}{c}{Correlation with Output} \\
\cmidrule(lr){4-7}\cmidrule(lr){9-13}          & $\sigma_y$     & $\sigma_{nx}$    & c     & x     & n     & z     &       & \multicolumn{1}{c}{c} & \multicolumn{1}{c}{x} & \multicolumn{1}{c}{nx} & \multicolumn{1}{c}{n} & \multicolumn{1}{c}{z}\\
\midrule
    Closed Economy & 1.80  &   & 0.35  & 3.58  & 0.58  & 0.50  &       & 0.94  & 0.80  &       & 0.93  & 0.90 \\
    Benchmark & 1.50  & 3.77  & 0.42  & 10.99 & 0.50  & 0.67  &       & 0.77  & 0.27  & 0.01  & 0.93  & 0.89 \\
    Transport Cost & 1.35  & 0.37  & 0.47  & 2.91  & 0.47  & 0.75  &       & 0.81  & 0.92  & 0.23  & 0.92  & 0.98 \\
    Autarky & 1.26  &       & 0.54  & 2.65  & 0.91  & 0.99  &       & 0.90  & 0.96  &       & 0.91  & 0.99 \\
  \bottomrule
    \end{tabular}
  \end{adjustbox}
\end{subtable}

\medskip

\begin{subtable}{\textwidth}
  \centering
  \caption{Comovements between the Two Countries}
    \begin{tabular}{lccccc}
    \toprule
          & \multicolumn{5}{c}{Cross-country Correlations} \\
      \cmidrule(lr){2-6}
            & y     & c     & x     & n     & z \\
    \midrule
    Benchmark & -.21 & .88  & -.94 & -.78 & .25 \\
    Transport Cost & -.05 & .89  & -.48 & -.70  & .25 \\
    Autarky & .08  & .56  & -.31 & -.51 & .25 \\
  \bottomrule
    \end{tabular}
\end{subtable}
\caption{Model Standard Deviations and Within- and Cross-Country Comovements}
\label{tab:modelresult}
\end{table}


\subsection{The Quantity Anomaly}

The simplest two-country RBC model, with free movement of goods and financial assets and no transport costs, differs from the data in several respects (\cref{tab:modelresult}: Benchmark). Investment and net exports are too volatile, consumption is too smooth, and output is negatively correlated across the two countries. The closed-economy model therefore reproduces these empirical features more successfully than its international counterpart. Most of the discrepancies can be reduced by changing the model's parameters or structure. For example, the model's net export and investment volatilities are initially seven and three times their empirical counterparts, respectively. Introducing quadratic transport costs brings them closer to the observed levels (\cref{tab:modelresult}: Transport Cost).

One discrepancy, however, is particularly robust. The model predicts that consumption is more highly correlated across countries than output, $\rho_{c,c^*} > \rho_{y,y^*}$, whereas the data show the opposite ranking, $\rho_{y,y^*} > \rho_{c,c^*}$. BKK call this discrepancy the quantity anomaly\footnote{BKK (1993) also identify a second anomaly: the standard deviation of the terms of trade is larger in the data than in the model. This paper focuses on the quantity anomaly.} because it persists across changes to both the model's parameters and its structure.

\subsection{Robustness of the Quantity Anomaly}

The baseline model produces a cross-country output correlation of -0.21. In BKK's data, by contrast, output correlations between the United States and the other ten countries are all positive, ranging from 0.38 for Austria to 0.76 for Canada. The model generates a negative correlation because a favorable technology shock raises the marginal productivity of labor and capital in the affected country. Higher returns attract productive resources from the other country, so output rises in one country at the expense of the other.

Consumption is perfectly correlated across countries when markets are complete, allowing international risk sharing, and preferences are additively separable between consumption and leisure (\cite{Scheinkman:1984aa}). BKK's benchmark model has complete markets but nonseparable preferences, so its consumption correlation is 0.88 rather than one. This is still much higher than the empirical correlation between the United States and Europe, which is 0.51.

To address the high consumption correlation and low output correlation, BKK eliminate international risk sharing altogether (\cref{tab:modelresult}: Autarky). Neither goods nor securities can be traded, leaving correlated technology shocks as the only link between the countries. Consumption correlation falls to 0.56, but output correlation remains low at 0.08, compared with 0.70 between the United States and Europe in the data. BKK suggest that the persistent consumption correlation under autarky may reflect the permanent income hypothesis. When a favorable technology shock occurs in one country, households in the other anticipate that its effects will spill over. They therefore increase current consumption and postpone investment.

Output and consumption correlations are closely linked in the model. Making productivity shocks more highly correlated raises cross-country output correlation but lowers consumption correlation. Different shock correlations can therefore reproduce either the observed output correlation or the observed consumption correlation, but not both simultaneously. The quantity anomaly concerns the relative magnitudes of these correlations, rather than either correlation in isolation.